\section{Who may seal}\label{sec:anchoring}

The sort so far says where authority is needed: wherever a grounding condition asks for an
absence. It says nothing yet about who holds that authority, how they came to hold it, or what
happens when they use it badly. A seal settles a question for everyone who accepts it, and whoever
has the standing to seal can also seal in bad faith. Signatures and validators offer no help with
this. A signature establishes that a key vouched for a record, and a validator establishes that the
record has the form its kind requires. A chair who closes the call for objections an hour early, a
clinician who signs within her rights for reasons she would never state and an administrator who
backdates a revocation all produce records that pass every check. This section asks what kind of
fact a charge of abuse is, where the rules that confer standing come from, how those rules change,
and how disputes about their use are settled.

\subsection{What a charge of abuse asks for}\label{sec:charge}

A charge that someone abused their standing mixes facts of different kinds, and the sort applies to
each part on its own.

Part of it is procedural. The chair gave notice of the vote on the fifth of the month; the call for
objections stayed open for three days; the endorsements came from members. Each of these asks only
that certain records exist, so each is monotone, and anyone holding the records can check it. A
procedure written into a frame turns a good share of what people mean by fair dealing into facts of
this kind.

Part of it asks for an absence. The chair had no undisclosed interest in the contract; no objection
was still open when the motion carried; the office-holder's grant had not been revoked when she
signed. These are the facts the sort singles out. Nobody can settle them from their own records, and
each needs a seal from someone with the standing to declare a set of records complete: a register of
interests closed on a date, a chair's call for objections, a revocation list signed as of a time.

The rest asks for something no record grounds: what the chair intended, whether the clinician's
reason was the one she gave, whether a member was pressed to vote outside every channel the records
cover. The camera that signed a scene it never photographed belongs with these. A signature shows
who stands behind a record and leaves open whether what the record reports took place. A body with
the standing to judge such facts can judge them, and its judgment then becomes a record like any
other, though no computation over records could have produced it.

Bad faith is therefore partly an absence fact and partly a fact about minds. The part that leaves
records can be made monotone by procedure or settled by a seal. The part that leaves none falls to
judgment, and a design of records can make that judgment, once given, attributable to whoever gave
it and open to appeal.

\subsection{Secondary rules as anchoring}\label{sec:secondary}

\citet{epstein2015} separates grounding, what makes a particular social fact hold, from anchoring,
what puts the frame principles in place. The sort applies to grounding conditions under a fixed
frame. For the frame itself, Hart's account of law supplies the vocabulary.

\citet[ch.~5]{hart2012} imagines a community governed only by primary rules, rules that require or
forbid conduct, and names the defects such a regime would suffer. It is uncertain, since the group
has no agreed way to identify which rules are its rules. It is static, since nothing provides for
changing the rules on purpose. And it is inefficient, since disputes over whether a rule was broken
have no one with the authority to settle them, and the pressure to comply stays diffuse. Hart pairs
each defect with a kind of secondary rule that remedies it: a rule of recognition, rules of change
and rules of adjudication.

A collaboration over signed records that has nothing beyond the records and each observer's trust
suffers all of Hart's defects. Which keys count as members is settled by each observer inside its
own fiber of trust \citep{karna2026}, so the group has no rule of recognition it can point to.
Nothing says how its frame changes, and a dispute over whether an absence held has nobody to seal
it. The objection from bad faith, in its strongest form, is aimed at a regime like this one, and
against such a regime it succeeds.

The remedy is to write the secondary rules down as records, co-signed by the parties they bind.
\begin{itemize}[leftmargin=1.5em,itemsep=2pt]
  \item A rule of recognition says which records count for the group: the founding record that
    names the members' keys and the group's constitutive rules, such as the committee's rule that a
    proposal goes to a vote once three members have signed endorsements of it.
  \item A rule of change says how the frame may change. A change of frame asks that no party still
    holds the old terms, which makes it absence-sensitive (the corollary to the second objection
    below), so a rule of change names the seal that settles each change.
  \item A rule of adjudication names who has the standing to seal a disputed absence, and the
    procedure they follow.
\end{itemize}
Each of these is a record its parties co-sign, so the frame is anchored in Epstein's sense, by the
collective acceptance of those it binds, and anyone bound by it can inspect that acceptance.
\citet{schneider2024} calls the default of online platforms \emph{implicit feudalism}, a bias,
cultural and technical, toward communities run as fiefdoms, and \citet{schneider2021} observe that
the governance tools online communities are given typically assign unchecked power to administrators
and moderators. A platform keeps its secondary rules in its operator's control of the servers, and
here the same rules are records signed by the people they bind.

\subsection{Frames in layers}\label{sec:layers}

Frames come in layers, and so do the vocabularies they are written in. Each person keeps an
idiolect, the olog of what their own records mean to them. A group, such as a family, a clinic or a
committee, keeps a sociolect, the ontolog its members have co-signed. A nation keeps a national
dialect, the legal and administrative vocabulary its institutions use. Rules, goals, targets and
policies sit at the same layers as the vocabularies: a patient's own targets and the plan she wrote
down for herself; a clinic's policies and a committee's rules; a nation's laws and the indicators it
reports.

The system these layers make up is built by the people in it, through consensus and unanimity. A
frame binds only the parties who co-signed it, and a procedure inside a frame, such as a quorum, a
majority vote or a panel, binds because everyone it binds co-signed the procedure first. A majority
vote then binds only those who agreed to be bound by majority votes. Unanimity across a nation is
out of reach, and unanimity among the parties a given rule binds comes within reach when the rule
lives at the smallest layer that contains all of them. Each layer is anchored by the collective
acceptance of its own members, and a rule that concerns one clinic stays with that clinic.

Groups customize the system by adding to it. A group adds the concepts of its sociolect, its own
rules, targets and policies, and the views its members read through, each co-signed by the members
it binds. A customization of this kind is monotone: it leaves in force everything the members and
the layers around them already hold, so it needs no coordination beyond the co-signatures of those
it binds. A change that takes something away asks for more. Narrowing a right a member holds, setting aside a rule of a wider layer, or rewriting a settled term asks that nobody still
relies on what is removed, which is an absence. Within its own frame a group can make such a change
only by the co-signature of every party it binds, the way the parties to an agreement amend it. A
change that takes something from a wider layer goes the way regulation changes, through the legal
pathways for adjusting it, where a body with the standing to seal decides.

Layers meet by translation. A term crosses from one vocabulary into another only when both sides
have signed that the two terms mean the same, and the translation carries a round-trip law:
exporting a record into the other vocabulary and importing it back returns the record unchanged.
Where a mapping is many-to-one, a lens keeps the part the wider vocabulary drops, so the round trip
still holds. A sociolect never overwrites the idiolects of its members, and a person's term reaches
the national vocabulary through the group's, with each step co-signed by the parties on both sides.

National law binds people who never co-signed it, and no frame can make it consensual. What a frame can do is record the point of contact: a group
co-signs a record that it takes up a national rule, and where a national rule clashes with a local
one, the clash sits on the record and goes up the chain of appeal that ends outside the
collaboration (\secref{sec:disputes}). Recognition then runs in both directions. The state recognizes
the group's right to organize, which is one of Ostrom's principles, and the group records which of
the state's rules it has taken up.

\subsection{The contractualist test, in tiers}\label{sec:tiers}

The frame we have defended elsewhere for collaborations of this kind holds rules to a
contractualist test: a rule that changes what others may do stands only if every party it binds has
co-signed it.\footnote{The test sits beside a second condition, that no party can write to another's
log, so a party bound by nothing it did not sign also keeps everything it did sign.} The test follows
\citet{scanlon1998}, on whose account an act is wrong when principles that no one could reasonably
reject, as a basis for informed and unforced general agreement, would disallow it. The co-signatures
are the record that nobody bound by the rule rejected it.

Applied to every decision, the test gives every party a veto, and a veto is a power that can be used
in bad faith too. A member who blocks every proposal until a side demand is met abuses her standing
as plainly as a chair who closes objections early. Scanlon's qualifier carries the weight: what
counts is reasonable rejection, and a rejection can be checked for reasons only if it states them.
Projects that run on consensus have reached the same requirement in practice. At the Apache Software
Foundation a negative vote from a qualified voter stops a proposed code change, and a veto without a
justification is invalid and has no weight \citep{apache-voting}. The contractualist test can take
the same form here: a refusal to co-sign stands as a rejection when it is a signed record carrying
its reason, and the reason is itself open to dispute under the rules of adjudication.

Even with reasoned vetoes, unanimity suits some decisions and not others. We propose tiers, each
fixed in the frame by the unanimous co-signature of its parties at the outset.
\begin{enumerate}[leftmargin=1.5em,itemsep=2pt]
  \item The constitutive rules, which make up the frame itself. A change needs the co-signature of
    every party the rule binds, and vetoes count with their reasons, so the contractualist test
    applies in full to the rules that define the group.
  \item Offices. The standing to seal is held on a grant that expires and is renewed by
    co-signature. A quorum named in the frame may suspend an office-holder at once; removal needs a
    higher threshold and a decision by a panel under the rules of adjudication.
  \item What concerns one person alone. A person's access to their own records, their leaving with
    their own history and their right to appeal are fixed in the frame as matters that need that
    person's own signature, and no quorum reaches them.
\end{enumerate}
Projects that govern themselves in the open set their thresholds in the same order. Debian's
developers may recall the project leader by a general resolution, overrule the technical committee
by a two-to-one majority, and amend their constitution only by three to one \citep{debian2022}.
Python's core team may remove a member of the steering council, or the whole council, by a vote of
no confidence that succeeds when at least two thirds of those voting express a lack of confidence
\citep{pep13}. In both, changing the frame costs more than using it, and removing an office-holder
costs more than the everyday decisions the office makes.

A recall under the second tier keeps faith with the joint commitment that created the office. On
Gilbert's account a joint commitment can be ended only by its parties together
\citep{gilbert1990,gilbert2013}, and here the parties used that power at the outset, when they
co-signed a frame that lets a named quorum suspend an office. A suspension follows the terms of the
commitment it suspends. The third tier gives Gilbert's point its widest reach: for what concerns one
person alone, that person is always among the parties who must sign.

\subsection{Who seals for the sealers}\label{sec:sealers}

Offices and panels hold the standing to seal, so their own make-up has to resist the members who
sit on them. A member with elevated standing can poison what the body seals from inside: by sealing
alone in the body's name, by backdating a record the body relies on, or by rotating out the members
who would refuse. Separating the body that suspends from the body that judges does not reach this,
since each body can be poisoned in the same way.

A governing body can be built as a composite collaborator, decentralized inside and presenting a
single facade outside. Inside, each member keeps their own log and key and works under the same
discipline as everyone else, adding records freely and sealing only where an absence matters.
Outside, the body acts through one facade, a signature that readers check without needing to know
how the body is wired. In the model of \citet{karna2026}, parties wired together present one
interface to the rest of a collaboration, and a governing body is one more composite of that kind,
standing as a facade over the layer it governs.

The facade's signature needs a quorum of the members. It can be a threshold signature, which
verifies as a single key once enough members have contributed their shares \citep{rfc9591}, or a
list of co-signatures under the record that founded the body. Either way each member's part stays
in the body's inner log, so the attribution that accountability depends on survives behind the
facade. Changing who sits on the body changes its frame and falls under the first tier, which keeps
a captured majority from rotating the honest members out on its own authority. Members hold
expiring grants, and the facade's key is renewed at each epoch, so a stolen or poisoned share loses
its value when the epoch ends. The facade that suspends and the facade that judges have disjoint
members.

A quorum raises the price of capture and leaves capture possible: enough colluding members can
still seal in the body's name. The regress ends at the founding co-signature that anchored the body
and at the institutions outside the collaboration where the chain of appeal ends.

\subsection{What a suspension costs}\label{sec:suspension}

A suspension is a revocation, and revocation asks for an absence: an office-holder's signature
counts only if no suspension of it came first. A suspension therefore takes effect for each reader
when that reader receives it, and an act the office-holder signed in the meantime is concurrent
with the suspension from that reader's point of view. No arrangement of signatures makes a
suspension take effect everywhere at once. The delay is part of what absence costs, and a frame can
only state it: a window within which a suspension is honored, paid for by fresh checkpoints that
readers fetch while the window lasts.

Concurrent revocation also opens a way to abuse standing, and builders of local-first software have
met it as the problem of duelling administrators. Administrators with equal rights revoke each
other concurrently, and the merged state does not say who prevails. A rule that breaks the tie by
seniority can be gamed: the senior administrator can backdate a revocation so that it appears
concurrent with the junior's, and so roll the junior's revocation back. \citet{dougal2026} argues
that a Byzantine administrator can exploit concurrency in this way to sway the outcome, so that an
arbiter outside the group's merge has to order concurrent events, and proposes optional epoch events
that give a bounded total order within each epoch. In the terms of this paper an epoch event does
the work of a seal: it declares which revocations count as earlier, and from then on their order is
a fact grounded in a record. An arbiter held by one party would be the elevated member of
\secref{sec:sealers}, and an arbiter built as a facade over a quorum orders the epoch without any
single member holding that power.\footnote{Other designs choose a different frame principle in
advance. A strong-removal rule lets a removal prevail over every concurrent act of the member
removed \citep{p2panda-auth}, which settles the duel without an arbiter by giving every removal the
same precedence, and moves the risk of abuse to whoever can remove.}

The case shows what a seal can do about bad faith and where it stops. The epoch event never detects
that the senior administrator backdated anything. A device's timestamp is a claim the device makes
about itself, and the records cannot ground its truth. What the seal takes away is the advantage:
once a witness the parties named in advance has fixed the order, a later claim about an earlier time
gains nothing. Much of the design of secondary rules has this shape, leaving intent to judgment and
taking the profit out of the abusive move.

\subsection{Disputes, sanctions and nested appeal}\label{sec:disputes}

\citet{ostrom1990} found that institutions which governed common resources for generations share a
set of design principles, among them monitors who are accountable to the users or are users
themselves, sanctions graduated by the seriousness and context of an offense, conflict-resolution
mechanisms that are cheap and close at hand, recognition by outside authorities of the users' right
to devise their own institutions, and, in larger systems, governance organized in nested layers. A
meta-analysis of the studies that tested the principles afterwards found them well supported
\citep{cox2010}. Each maps onto a design of records without strain.

Monitoring comes almost free, since every act is a signed record and its signer is known. What
remains is that the people an act concerns can see it, which is a question of who receives which
records. Sanctions can be graduated as records of increasing weight: a concern noted on the record,
a suspended grant, a recall, a removal, each sealed by whoever the frame gives standing for it. A
dispute is a record that opens a procedure, and the procedure has the shape of a path of arrows:
filed, answered, decided, appealed, each step a signed record that answers the one before, with the
decision a seal on what the dispute considered.\footnote{In \citeauthor{hirschman1970}'s
(\citeyear{hirschman1970}) terms, refusing a party's records inside one's own fiber of trust is exit,
and a dispute is voice. The design keeps exit open for everyone at no cost, and the rules of
adjudication are what make voice worth using for those, such as patients and employees, who cannot
easily leave.}

Nesting is where a design of records meets the institutions around it, and the layers of
\secref{sec:layers} give the nesting its order. The Internet Engineering Task Force gives a working
example of nested appeal among peers. Someone who disagrees with a working group's decision takes it
first to the group's chairs, then to the responsible area director, then to the steering group and
finally to the architecture board, whose decision is final \citep[\S6.5.1]{rfc2026}, and every
activity in the process is kept in a publicly accessible record \citep[\S8]{rfc2026}. At the top of
any such chain sits a body outside the collaboration, such as a professional council or a court.
Searle's reply to Smith's cases of contested property, that problems of that kind are addressed by
legal or political means \citep{smithsearle2003}, puts it there, and Ostrom's principle of outside
recognition says why a chain of appeal needs to end in an institution that recognizes the group's
own rules.

Who should sit on a panel is a question the evidence leaves open. In a survey experiment with
Facebook users, \citet{pan2022} found expert panels perceived as more legitimate than paid
contractors, algorithms or juries of users, and found that agreeing with a decision mattered more to
its perceived legitimacy than the process that produced it. A frame can choose either kind of panel,
and it gains from making the choice in advance and on the record, before any particular outcome is
at stake.

\subsection{Stake as a record}\label{sec:stake}

A familiar answer to abuse is to make office-holders stake something they can lose, money or
reputation. A design of records can carry reputation without a new institution, since a person's
signed history is already in the records: the decisions they sealed, the disputes decided against
them, the suspensions they served. We keep that history as a record and leave it unscored. Each
observer reads it inside its own fiber, as trust is read everywhere else in the design, and nobody
publishes a number about a person. A published score would rank people for every observer at once,
which would bring back the single authority the design keeps out, and would hand anyone set on abuse
a single target to game. Money and licenses stay with the institutions that already hold them under
law: a clinician's license with the council that grants it, a bond with a regulated custodian. The
records give those institutions what they need in order to act, an attributable history of what the
office-holder signed.

\subsection{The imagined clinic}\label{sec:clinic}

In the imagined clinic, a clinician's standing to sign prescriptions for the clinic rests on a grant
from the clinic that expires and is renewed. The clinic's frame names who may grant and suspend it,
for instance its medical director together with one other clinician, acting as one facade. A
patient who believes a clinician has abused that standing can move at every tier. Inside her own
fiber she can stop treating the clinician's new records as ones she acts on, which needs nobody's
agreement, and she keeps her own records whatever happens, since they concern her alone. She can
file a dispute, a signed record that opens the clinic's procedure. The clinic's panel can suspend
the grant as a sanction, and the suspension reaches the clinic's devices within the window its frame
states, under a witness's seal, so that no prescription signed afterwards can be backdated into the
gap. If the clinic's own procedure fails her, her last appeal goes to the council that licensed the
clinician, which holds the stake that matters most and can judge what the records cannot, namely why
the clinician signed.

\subsection{Against the central alternative}\label{sec:central}

The comparison with a central platform turns on where the seals are and who can see them. A platform
seals every fact, monotone or not, and its sealing leaves records that only the operator reads.
Abuse within the operator's standing, or within the standing of the moderators it appoints, stays
out of sight of those it affects unless the operator chooses to show it. A collaboration that keeps
its secondary rules as records seals only where absence matters, every seal is a signed record
naming its signers, and every office rests on a grant that expires. Abuse there becomes attributable
and contestable. Accountability needs both, and it also needs a body that hears disputes, since
attribution with nobody to hear the dispute leaves the injured party holding a well-documented
grievance and nowhere to take it. A collaboration with records and no secondary rules can therefore
end up worse than a platform, which at least has an operator to complain to. With its secondary
rules written as records and anchored by those they bind, the collaboration offers what a platform
withholds: seals in view, offices that lapse, and appeals that end in institutions outside it.

Collectives abuse standing too, and the layers and tiers are meant for them as much as for
office-holders. A collective binds only the parties who signed; changing its frame needs every bound
party's co-signature, with reasoned vetoes; its customizations add to the wider layers and take
nothing from them except through the legal pathways; suspensions are quick and reversible while
removals are slow; the body that suspends differs from the one that judges; what concerns one person
alone lies beyond every quorum; and the chain of appeal ends outside the collaboration. Abuse stays
possible under these rules, and it moves slowly, in view, against people who keep a remedy. Who
holds a seal remains a political question, and the sort shows only where that question has to be
asked.
